
\section{Reflektion över kunskaper och färdigheter}
\label{sec:reflektion}

LIA 2 har gett mig en förståelse för produktionssystem som jag inte hade kunnat få enbart genom
kurser. Jag har inte bara lärt mig nya verktyg, utan också utvecklat mitt sätt att tänka,
planera, samarbeta och ta ansvar.

\subsection*{Miljön och förutsättningarna}

Insighta gav mig tillgång till GCP-projekt, GKE-kluster, Terraform-kodbas och GitHub-repot.
Jag deltog i stand-ups, tekniska möten och code reviews precis som övriga teammedlemmar.
Det var alltså inte övningsuppgifter, utan riktigt arbete med riktiga konsekvenser.

Kommunikationen skedde främst via Slack, GitHub pull requests och möten. Eftersom teamet är
internationellt blev skriftlig kommunikation extra viktig. Jag märkte snabbt att en tydlig pull
request sparar mycket tid för hela gruppen.

\subsection{Feedback jag fått och vad den sagt mig}
\label{subsec:feedback}

\subsubsection{Tekniska kunskaper}
\label{subsubsec:kunskaper_feedback}

Jag fick positiv feedback på att jag tog mig an ett komplext projekt och drev det fram till en
fungerande lösning. Teamet lyfte särskilt att jag förstod sambandet mellan IAM, RBAC och
Workload Identity och att jag valde en säker lösning i stället för den snabbaste.

Jag tog hänsyn till den feedbacken genom att fortsätta arbeta på samma sätt i resten av
perioden. När jag ställdes inför nya tekniska val försökte jag därför inte bara lösa problemet
snabbt, utan också tänka på säkerhet, tydlighet och långsiktig hållbarhet.

\subsubsection{Planering och prioritering}
\label{subsubsec:omdome_feedback}

Jag fick höra att jag hade blivit bättre på att avgränsa och prioritera mitt arbete. Mina pull
requests var tydligare avgränsade och lättare att granska än tidigare.

Jag tog till mig den feedbacken genom att fortsätta dela upp stora uppgifter i mindre delar.
Det hjälpte mig att hålla bättre kontroll över arbetet och gjorde det lättare att snabbt upptäcka
var ett fel hade uppstått. Att jag under perioden skapade 21 pull requests i detta repo visar
också att jag faktiskt arbetade på det sättet i praktiken. Jag försökte göra varje ändring liten
nog för att vara tydlig, men tillräckligt stor för att föra arbetet framåt.

\subsubsection{Kommunikation och samarbete}
\label{subsubsec:samarbete_feedback}

Jag fick positiv feedback på hur jag kommunicerade i Slack och i mina pull requests. Teamet
uppskattade att jag beskrev vad jag gjort, varför jag gjort det och hur ändringen kunde testas.
Jag fick också beröm för att jag bad om hjälp i rätt tid när jag fastnade.

Jag tog hänsyn till detta genom att bli mer medveten om hur jag kommunicerade. Jag försökte
skriva tydligare, ge bättre statusuppdateringar och tidigare lyfta hinder som kunde påverka
arbetet.

\subsubsection{Initiativförmåga och ansvar}
\label{subsubsec:prestation_feedback}

Teamet uppskattade att jag tog initiativ även utanför de uppgifter jag först fått. Ett exempel var
att jag såg problem i deploymentflödet som inte uttryckligen stod i uppgiften, men som ändå
behövde förbättras.

Jag tog till mig detta genom att börja se mer till helheten. Jag försökte inte bara lösa min
närmaste uppgift, utan också fundera på om det fanns något runt omkring som kunde göras
bättre för teamet.

\subsection{Hur jag har utvecklat mina kunskaper och färdigheter}
\label{subsec:utveckling}

\subsubsection{Tekniska kunskaper}
\label{subsubsec:kunskaper_utveckling}

Den största tekniska utvecklingen under LIA 2 har varit min förståelse för säkerhet,
identitet och infrastruktur som kod. I skolan lärde jag mig vad IAM och RBAC är. Under LIA 2
lärde jag mig hur de faktiskt samverkar i ett riktigt system.

Jag utvecklade också min förståelse för Terraform. Det handlar inte bara om att skapa resurser,
utan också om att skriva tydlig kod som går att underhålla och vidareutveckla.

\subsubsection{Planering, prioritering och omdöme}
\label{subsubsec:omdome_utveckling}

Jag har blivit bättre på att planera och prioritera genom att dela upp stora problem i mindre,
testbara steg. Det gjorde att jag snabbare kunde hitta fel och rätta dem innan de växte till
större problem.

Jag märkte också att denna förmåga gjorde mig lugnare i arbetet. När ett problem känns stort
blir det lättare att hantera om man bryter ned det i mindre delar.

\subsubsection{Kommunikation och samarbete}
\label{subsubsec:samarbete_utveckling}

Jag har blivit bättre på att skriva tydliga tekniska förklaringar. Tidigare tänkte jag mest på att
koden skulle fungera. Nu tänker jag också på att andra ska kunna förstå varför jag valt en viss
lösning.

Det är särskilt viktigt i ett distribuerat team där mycket kommunikation sker skriftligt.

\subsubsection{Ansvar och initiativförmåga}
\label{subsubsec:prestation_utveckling}

Den tydligaste personliga utvecklingen under LIA 2 är att jag har tagit större ansvar. Jag ser
nu tydligare att mitt arbete påverkar helheten. Om jag upptäcker ett problem räcker det inte att
tänka att det ligger utanför min egen uppgift. Jag behöver åtminstone lyfta det och ibland även
ta initiativ till att lösa det.